\documentclass[conference]{IEEEtran}

\usepackage{amsmath}
\usepackage{booktabs}
\usepackage{graphicx}
\usepackage{url}

\title{Cluster-Local Streaming Merge-Update Offload for Online Softmax State Updates on Spatz}

\author{
  \IEEEauthorblockN{Wenxuan Tang}
  \IEEEauthorblockA{Local Spatz Prototype Study}
}

\begin{document}

\maketitle

\begin{abstract}
Online softmax-style state updates maintain short-lived scalar and vector
state across streaming tiles~\cite{online-softmax2018}. On a compact vector
cluster, these updates can be
awkward for a tensor- or vector-centric execution path: the work mixes scalar
state, vector payloads, fine-grained memory traffic, and a strict update order.
This paper presents a cluster-local streaming merge-update offload prototype
for Spatz. The design adds a memory-mapped control interface and one TCDM
master port, leaving the Spatz ISA, decoder, vector functional units, and
register files unchanged. The current prototype intentionally implements a
restricted merge semantics: zero-length old or tile states and equal-weight
same-exponent merges. Under this restricted semantics, a Verilator
microbenchmark shows that the engine is slower for tiny workloads because of
fixed MMIO and polling overheads, but reaches 2.40$\times$ speedup and a
58.3\% cycle reduction at $N=16,D=64$. A finer break-even sweep places the
threshold between roughly 64 and 96 vector elements for the equal-weight case.
The result supports cluster-local offload as a promising mechanism for
state-update-heavy online softmax kernels, while making clear that complete
mixed-scalar online softmax and end-to-end attention acceleration require a
future datapath with exponential, scaling, and normalization support.
\end{abstract}

\section{Introduction}

Attention implementations often rely on online softmax formulations to merge
partial summaries without materializing every intermediate score. Each merge
updates a scalar maximum $m$, a scalar normalizer $l$, and an output vector
$O$. The operation has a small amount of arithmetic per element, a strict state
dependency per row, and frequent access to short-lived intermediate state. This
shape does not naturally match a path optimized for large regular vector or
matrix kernels.

This work studies whether a small cluster-local offload block can handle this
state-update pattern more efficiently than scalar software on the cluster core.
The implemented prototype is deliberately conservative. It does not introduce
new instructions and does not modify Spatz pipeline files. Instead, software
programs a memory-mapped merge engine, the engine reads and writes TCDM data
through its own master port, and software observes completion through status
bits.

The paper makes three contributions:

\begin{itemize}
  \item It integrates a streaming merge-update engine into the Spatz cluster
        through MMIO control and a TCDM master port, without changing the ISA or
        vector pipeline.
  \item It defines and validates a restricted online-softmax merge prototype,
        including correct outputs for supported cases and explicit error
        reporting for invalid or unsupported inputs.
  \item It evaluates CPU scalar execution against the offload path in
        Verilator, reporting cycles, break-even behavior, and TCDM access and
        congestion counters.
\end{itemize}

The scope is also explicit. This is not yet a full online softmax datapath and
does not claim end-to-end attention speedup. The missing pieces are
exponential evaluation, weighted scaling, and normalization division for
general mixed-scalar inputs.

\section{Background and Motivation}

\subsection{Online Softmax Merge State}

For a complete online softmax merge, two summaries are combined as:

\begin{align}
m_{new} &= \max(m_{old}, m_{tile}), \\
l_{new} &= l_{old}e^{m_{old}-m_{new}} +
           l_{tile}e^{m_{tile}-m_{new}}, \\
O_{new} &= O_{old}
          \frac{l_{old}e^{m_{old}-m_{new}}}{l_{new}} +
          O_{tile}
          \frac{l_{tile}e^{m_{tile}-m_{new}}}{l_{new}}.
\end{align}

The current RTL prototype implements only a subset of this equation:
\texttt{l\_old=0}, \texttt{l\_tile=0}, and constructed equal-weight cases where
\texttt{m\_old==m\_tile} and \texttt{l\_old==l\_tile}. Other legal
mixed-scalar inputs are detected and reported as unsupported errors rather
than producing silent approximations.

\subsection{Spatz and TCDM}

Spatz is a compact vector-capable cluster system built around compact
RISC-V-based vector units and Snitch scalar cores~\cite{spatz2023,snitch2020}.
It provides tightly coupled data memory (TCDM) close to the cores and
accelerators. The online merge prototype
uses this memory hierarchy directly: all input and output buffers reside in
TCDM, and the engine issues its own TCDM requests. This placement keeps the
offload local to the cluster and avoids changing the external memory system.

\subsection{Why Offload This Pattern}

The merge-update operation is small, ordered, and stateful. A scalar software
implementation pays instruction overhead for every row and vector element. A
streaming engine can amortize setup overhead over many elements, use a fixed
state machine for memory scheduling, and avoid repeatedly re-executing the same
control sequence in software. The key question is whether the fixed MMIO and
polling overheads can be amortized for practical problem sizes.

\section{Design}

\subsection{Software-Visible Interface}

Software configures the engine using memory-mapped registers. Table
\ref{tab:mmio} summarizes the interface.

\begin{table}[t]
\centering
\caption{Memory-mapped merge engine interface.}
\label{tab:mmio}
\begin{tabular}{ll}
\toprule
Register & Purpose \\
\midrule
\texttt{MERGE\_SRC\_M\_OLD} & Base address of old $m$ array \\
\texttt{MERGE\_SRC\_L\_OLD} & Base address of old $l$ array \\
\texttt{MERGE\_SRC\_O\_OLD} & Base address of old $O$ matrix \\
\texttt{MERGE\_SRC\_M\_TILE} & Base address of tile $m$ array \\
\texttt{MERGE\_SRC\_L\_TILE} & Base address of tile $l$ array \\
\texttt{MERGE\_SRC\_O\_TILE} & Base address of tile $O$ matrix \\
\texttt{MERGE\_DST\_M} & Base address of output $m$ array \\
\texttt{MERGE\_DST\_L} & Base address of output $l$ array \\
\texttt{MERGE\_DST\_O} & Base address of output $O$ matrix \\
\texttt{MERGE\_N} & Number of rows \\
\texttt{MERGE\_D} & Vector dimension per row \\
\texttt{MERGE\_STRIDE} & Row stride in bytes; zero means packed \\
\texttt{MERGE\_CTRL} & \texttt{start} and \texttt{clear\_done} pulses \\
\texttt{MERGE\_STATUS} & \texttt{busy}, \texttt{done}, and \texttt{error} bits \\
\bottomrule
\end{tabular}
\end{table}

The software wrapper writes source and destination offsets, dimensions, and
stride, then pulses \texttt{start}. Completion is detected by polling
\texttt{MERGE\_STATUS}. Invalid dimensions, misaligned addresses, and
unsupported mixed-scalar cases set \texttt{error}.

\subsection{Engine Microarchitecture}

The engine is a streaming state machine. On \texttt{start}, it validates the
configuration, reads scalar state for one row, computes restricted scalar
merge results, stores scalar outputs, streams the $D$ vector elements, and then
advances to the next row. The conceptual states are:
\texttt{IDLE}, \texttt{LOAD\_SCALAR}, \texttt{COMPUTE\_SCALAR},
\texttt{STORE\_SCALAR}, \texttt{UPDATE\_VECTOR}, \texttt{DONE}, and
\texttt{ERROR}, as summarized in Figure~\ref{fig:engine-flow}.

\begin{figure}[t]
\centering
\includegraphics[width=\linewidth]{../data_process/attnres/pic/online_softmax_merge_engine_flow.png}
\caption{Streaming merge-update engine control flow.}
\label{fig:engine-flow}
\end{figure}

The current datapath intentionally uses simple special-case floating-point
operations instead of a complete exponential and division pipeline. This keeps
the prototype focused on integration, memory scheduling, and offload overhead.

\subsection{Cluster Integration}

The engine is instantiated in \texttt{spatz\_cluster.sv}. A new TCDM request
and response channel is added to the narrow TCDM interconnect alongside core
and AXI-to-TCDM traffic. The request metadata marks engine traffic as non-core
traffic, allowing the existing performance counter path to include the new
port. The source file is added to the build manifest so Verilator builds the
engine as part of the cluster. Figure~\ref{fig:cluster-integration} shows the
integration point.

\begin{figure}[t]
\centering
\includegraphics[width=\linewidth]{../data_process/attnres/pic/online_softmax_merge_cluster_integration.png}
\caption{Cluster-local integration of the merge engine.}
\label{fig:cluster-integration}
\end{figure}

\subsection{Supported and Unsupported Semantics}

Table \ref{tab:semantics} records the tested input classes. The distinction is
important because the prototype is a merge-update offload experiment, not a
complete online softmax implementation.

\begin{table}[t]
\centering
\caption{Supported and unsupported input semantics.}
\label{tab:semantics}
\begin{tabular}{lll}
\toprule
Class & Condition & Behavior \\
\midrule
Case 0 & \texttt{l\_tile=0} & Select old state \\
Case 1 & \texttt{l\_old=0} & Select tile state \\
Case 2 & Equal $m$ and equal $l$ & Equal-weight merge \\
Case 3/4 & Equal $m$ and small equal $l$ & Equal-weight merge \\
Invalid & Zero size or misaligned config & Error \\
Unsupported & General mixed scalar state & Error \\
\bottomrule
\end{tabular}
\end{table}

\section{Implementation}

The implementation touches both hardware and software and is documented in the
local experiment notes~\cite{spatz-local-docs}. The main RTL module is
\texttt{hw/ip/online\_merge/src/online\_merge\_update\_engine.sv}. The cluster
top-level integration is in \texttt{hw/system/spatz\_cluster/src/spatz\_cluster.sv}.
The peripheral register description adds the \texttt{MERGE\_*} registers and
regenerates the SystemVerilog and C headers. The runtime-visible definitions
are exposed through \texttt{sw/snRuntime/include/spatz\_cluster\_peripheral.h}.

The benchmark is \texttt{sw/spatzBenchmarks/online-softmax-merge/main.c}. It
allocates TCDM buffers, initializes supported cases, computes a CPU reference,
starts the engine, checks outputs, records cycle counts, and reports TCDM
accessed and congested counters. It also runs invalid-configuration tests and
unsupported mixed-scalar probes.

\section{Evaluation}

\subsection{Experimental Setup}

All results in this paper use the Verilator target from the Spatz cluster build
flow. The software was rebuilt with:

\begin{verbatim}
make -C hw/system/spatz_cluster sw.vlt
\end{verbatim}

The benchmark was executed three times with:

\begin{verbatim}
ctest -R online-softmax-merge -V
\end{verbatim}

The three runs passed and produced identical cycle and TCDM counter values for
all valid cases. The CTest wall times were 229.02 s, 232.16 s, and 219.20 s,
but wall time is not used as an architectural performance metric.

\subsection{Correctness and Error Handling}

For all supported cases, the engine output matches the CPU reference within
the benchmark tolerance. The benchmark also checks error paths:
\texttt{N=0}, \texttt{D=0}, misaligned address, misaligned stride, mixed $m$,
unequal $l$, and a full-reference mixed-scalar probe. Unsupported legal
mixed-scalar inputs report \texttt{MERGE\_STATUS.error}; this is expected for
the restricted prototype.

\subsection{Cycle Results}

Figure \ref{fig:cycles} compares CPU scalar cycles and engine cycles. Tiny
inputs are slower on the engine because the fixed setup and polling overheads
dominate. As the vector work grows, the engine amortizes the setup cost.
Table~\ref{tab:perf} lists representative points from the same run.

\begin{figure}[t]
\centering
\includegraphics[width=\linewidth]{../data_process/attnres/pic/online_softmax_merge_bypass_cycles.png}
\caption{CPU scalar path versus merge engine path cycles.}
\label{fig:cycles}
\end{figure}

\begin{table}[t]
\centering
\caption{Selected performance points.}
\label{tab:perf}
\begin{tabular}{rrrrr}
\toprule
$N$ & $D$ & Case & CPU & Engine \\
\midrule
1 & 1 & 0 & 191 & 1216 \\
4 & 8 & 1 & 915 & 1486 \\
8 & 8 & 2 & 1704 & 1789 \\
8 & 12 & 2 & 2433 & 2011 \\
8 & 16 & 2 & 3114 & 2254 \\
8 & 32 & 2 & 5910 & 3284 \\
16 & 64 & 2 & 23033 & 9607 \\
\bottomrule
\end{tabular}
\end{table}

\subsection{Speedup and Break-Even}

Figure \ref{fig:speedup} reports speedup as CPU cycles divided by engine
cycles. The largest tested point, $N=16,D=64$, reaches 2.40$\times$ speedup.
The corresponding cycle reduction is:

\begin{equation}
\frac{23033 - 9607}{23033} = 58.3\%.
\end{equation}

\begin{figure}[t]
\centering
\includegraphics[width=\linewidth]{../data_process/attnres/pic/online_softmax_merge_bypass_speedup.png}
\caption{Speedup of the merge engine path over the CPU scalar path.}
\label{fig:speedup}
\end{figure}

The finer equal-weight sweep in Figure \ref{fig:breakeven} shows the
break-even region. $N=4,D=16$ and $N=8,D=8$ are still slightly below 1.0$\times$,
while $N=8,D=12$ reaches 1.21$\times$. The break-even point for this restricted
case is therefore between roughly 64 and 96 vector elements.

\begin{figure}[t]
\centering
\includegraphics[width=\linewidth]{../data_process/attnres/pic/online_softmax_merge_bypass_break_even.png}
\caption{Break-even sweep for restricted equal-weight case 2.}
\label{fig:breakeven}
\end{figure}

\subsection{TCDM Counters}

Figure \ref{fig:tcdm} shows the TCDM counters measured during engine execution.
The accessed counter scales with workload size, confirming that the engine
uses its own TCDM master path. Congestion is nonzero but much smaller than the
access count in this single-engine microbenchmark.

\begin{figure}[t]
\centering
\includegraphics[width=\linewidth]{../data_process/attnres/pic/online_softmax_merge_bypass_tcdm.png}
\caption{TCDM accessed and congested counters for the engine path.}
\label{fig:tcdm}
\end{figure}

\section{Limitations}

The main limitation is semantic: the RTL does not implement the full online
softmax equation. It lacks a general exponential datapath, weighted
multiplication for arbitrary scalar summaries, and division or reciprocal
normalization. Consequently, the reported speedups apply only to the supported
restricted cases.

The second limitation is workload scope. The evaluation is a microbenchmark
that isolates merge-update behavior. It does not prove end-to-end attention or
LLM inference speedup. AttnRes software-side traffic experiments are useful
motivation for state reuse, but they are not direct hardware performance
measurements for this engine.

\section{Related Work}

\textbf{Online softmax.}
Milakov and Gimelshein introduced online normalizer calculation for softmax,
showing that the normalizer can be maintained while reducing memory
accesses~\cite{online-softmax2018}. This paper uses the same broad motivation:
softmax state can be updated in a streaming fashion, and reducing unnecessary
memory traffic is valuable. Our prototype does not yet implement their full
generality; it isolates a restricted merge-update subset to study cluster-local
offload overheads.

\textbf{IO-aware attention.}
FlashAttention makes exact attention IO-aware by tiling the computation and
reducing traffic between GPU HBM and on-chip SRAM~\cite{flashattention2022}.
Our work is much narrower and targets a different platform. Rather than
restructuring a complete GPU attention kernel, it evaluates whether the
state-update part of an online-softmax-like flow is worth offloading near
cluster-local TCDM.

\textbf{Compact cluster processors.}
Snitch explores a tiny scalar core paired with floating-point execution and
minimal ISA extensions to improve energy efficiency for floating-point
workloads~\cite{snitch2020}. Spatz builds compact RISC-V vector units around
this style of cluster system~\cite{spatz2023}. The merge-update engine in this
paper follows a different integration point: it avoids ISA changes and uses a
memory-mapped, cluster-local TCDM master instead of extending the core or vector
pipeline.

\section{Conclusion}

This paper presented a cluster-local streaming merge-update offload prototype
for Spatz. The prototype demonstrates that a small MMIO-controlled TCDM master
engine can outperform scalar software for sufficiently large restricted
online-softmax state-update workloads, reaching 2.40$\times$ speedup and a
58.3\% cycle reduction at $N=16,D=64$. The break-even sweep shows that fixed
offload overheads are amortized between roughly 64 and 96 vector elements for
the equal-weight case. These results justify a next-stage design with a
complete online softmax datapath, while the current paper remains scoped to the
validated restricted prototype.

\bibliographystyle{IEEEtran}
\bibliography{common/bib/refs}

\end{document}
